%% ma: L12-B19-0034
%% lop: 12
%% chuong: 6
%% bai: 19
%% dang: 12.19.5
%% loai: TLN
%% muc: VD
%% dap_an: "0,03"
%% nguon: "(NBV)-10 ĐỀ KIỂM TRA HK2-TOÁN 12-DUNG CHUNG 3 SÁCH.pdf (H03, Đề số 3), Phần 3 Câu 1"
%% kien_thuc: "Bài 19 (công thức xác suất toàn phần và công thức Bayes)"
%% trang_thai: chinh-thuc
%% ghi_chu: "Đáp án gốc 0,03 khớp. Bỏ ảnh bác sĩ khám bệnh (ảnh trang trí)."
\begin{ex}
Trong một đợt kiểm tra sức khoẻ, có một loại bệnh $X$ mà tỉ lệ người mắc bệnh là $0{,}2\%$ và một loại xét nghiệm $Y$ mà ai mắc bệnh $X$ khi xét nghiệm $Y$ cũng có phản ứng dương tính. Tuy nhiên, có $6\%$ những người không bị bệnh $X$ lại có phản ứng dương tính với xét nghiệm $Y$. Chọn ngẫu nhiên $1$ người trong đợt kiểm tra sức khoẻ đó. Giả sử người đó có phản ứng dương tính với xét nghiệm $Y$. Xác suất người đó bị mắc bệnh $X$ là bao nhiêu (làm tròn kết quả đến hàng phần trăm)?
\shortans{0,03}
\loigiai{Gọi $A$: ``người được chọn mắc bệnh $X$'', $B$: ``người được chọn có phản ứng dương tính''. $P(A)=0{,}002$, $P\left(\overline A\right)=0{,}998$, $P(B\mid A)=1$, $P\left(B\mid\overline A\right)=0{,}06$.

Công thức Bayes: $P(A\mid B)=\dfrac{0{,}002\cdot1}{0{,}002\cdot1+0{,}998\cdot0{,}06}\approx0{,}0323\approx0{,}03$.}
\end{ex}
